\chapter{Общий родитель проводимости хопфовского цикла}
\label{ch:v10-hopf-cycle-conductance-parent}

\section{Связывание двух динамических скоростей}

Предыдущий гейт построил канонический марковский цикл в
$\mathcal K_{43}\otimes\mathcal H_{\rm KMS}$, но оставил его общую
проводимость $\kappa$ свободной. В проекте уже имеется вторая физически
типизированная скорость --- скорость рождения ячеек $\Gamma_B$, связанная с
частотой квантовых часов $\Omega$ посредством
\begin{equation}
 k_X:=\log\frac{1+2x}{1+x}=3\Delta\zeta,
 \qquad
 \frac{\Gamma_B}{\Omega}=k_X,
 \qquad x=e^{-S_{\rm vac}}>0.
 \label{eq:v10-conductance-growth-coupling}
\end{equation}
Проверим, может ли один положительный родитель отождествить скорость
прохождения через цикл со скоростью рождения геометрии.

Введём два безразмерных отношения
\begin{equation}
 r_B:=\frac{\Gamma_B}{\Omega},
 \qquad
 r_\kappa:=\frac{\kappa}{\Omega}.
\end{equation}
Минимальный общий родитель имеет вид
\begin{equation}
 \mathcal P_\kappa(r_B,r_\kappa)
 =\frac12(r_B-k_X)^2+\frac12(r_\kappa-r_B)^2.
 \label{eq:v10-conductance-common-parent}
\end{equation}
Он ограничен снизу и обращается в нуль только при
\begin{equation}
 r_B=r_\kappa=k_X.
 \label{eq:v10-conductance-selected-point}
\end{equation}
Его гессиан равен
\begin{equation}
 \operatorname{Hess}\mathcal P_\kappa
 =\begin{pmatrix}2&-1\\-1&1\end{pmatrix},
 \qquad \rank=2,
 \qquad \det=1,
 \label{eq:v10-conductance-parent-hessian}
\end{equation}
и имеет положительный спектр
$\{(3-\sqrt5)/2,(3+\sqrt5)/2\}$.

\section{Слепое совпадение тока и геометрического хода}

Стационарная точка родителя даёт
\begin{equation}
 \kappa=\Gamma_B=k_X\Omega.
 \label{eq:v10-conductance-birth-identification}
\end{equation}
Для тока на одном ребре хопфовского цикла отсюда следует
\begin{equation}
 \frac{J_{\rm edge}}{\Omega}
 =\frac{\kappa}{3\Omega}
 =\frac{k_X}{3}
 =\Delta\zeta.
 \label{eq:v10-conductance-blind-current}
\end{equation}
Это первое прямое совпадение между сквозным током цикла и геометрическим
ходом рождения ячейки: за один безразмерный такт часов ток через ребро равен
приращению логарифмического масштаба.

Производство энтропии также получает безразмерную нормировку:
\begin{equation}
 \frac{\sigma_\circlearrowright}{\Omega}
 =\frac{\kappa\log2}{\Omega}
 =3\Delta\zeta\log2.
 \label{eq:v10-conductance-blind-entropy}
\end{equation}

\section{Оставшаяся абсолютная орбита}

В логарифмических переменных $(\kappa,\Gamma_B,\Omega)$ два условия
$\Gamma_B/\Omega=k_X$ и $\kappa/\Gamma_B=1$ имеют матрицу
\begin{equation}
 C_\kappa=
 \begin{pmatrix}
 0&1&-1\\
 1&-1&0
 \end{pmatrix},
 \qquad \rank C_\kappa=2,
 \qquad \ker C_\kappa=\operatorname{span}\{(1,1,1)^T\}.
 \label{eq:v10-conductance-scale-map}
\end{equation}
Поэтому преобразование
\begin{equation}
 (\kappa,\Gamma_B,\Omega)\longmapsto
 c(\kappa,\Gamma_B,\Omega),
 \qquad t\longmapsto t/c
 \label{eq:v10-conductance-scale-orbit}
\end{equation}
сохраняет родитель и все полученные отношения. Общий родитель действительно
связывает две динамики, но не создаёт абсолютную секунду.

\begin{theorem}[Совместное относительное происхождение проводимости]
Положительный родитель $\mathcal P_\kappa$ единственным образом выбирает
$\kappa=\Gamma_B=k_X\Omega$. Поэтому стационарный ток хопфовского цикла и
геометрический ход удовлетворяют слепому равенству
$J_{\rm edge}/\Omega=\Delta\zeta$, а производство энтропии удовлетворяет
$\sigma_\circlearrowright/\Omega=3\Delta\zeta\log2$.
\end{theorem}

\begin{nogocriterion}[Общий относительный родитель не создаёт секунду]
Если родитель зависит только от отношений
$\Gamma_B/\Omega$, $\kappa/\Gamma_B$ и нормированной меры рождения, то
совместное масштабирование трёх частот остаётся точной нулевой модой.
Следовательно, равенство $\kappa=\Gamma_B$ является физически содержательной
связью двух процессов, но не выводом абсолютной проводимости.
\end{nogocriterion}

\begin{protocol}[Статус родителя проводимости]\begin{enumerate}\raggedright
\item положительный общий родитель: \textbf{построен};
\item проводимость относительно часов: \textbf{выбрана};
\item связь $\kappa=\Gamma_B$: \textbf{получена};
\item слепое равенство $J_{\rm edge}/\Omega=\Delta\zeta$: \textbf{получено};
\item архитектура: \textbf{$10/10$};
\item относительное происхождение: \textbf{$6/6$};
\item происхождение общего родителя: \textbf{$1/1$};
\item абсолютная проводимость и часы: \textbf{$0/2$};
\item следующий гейт: \textbf{аудит космологической постоянной как якоря
проводимости}.
\end{enumerate}\end{protocol}

Машинный сертификат сохранён в
\path{s2t_v10_cell_birth_four_volume_hopf_cycle_conductance_common_parent_origin_gate_results.json}.